\documentclass[aps,pra,twocolumn,superscriptaddress,longbibliography]{revtex4-2}

\usepackage{amsmath,amssymb,amsthm,graphicx,xcolor,bm,braket,dcolumn,booktabs}
\usepackage{hyperref}

\newtheorem{theorem}{Theorem}
\newtheorem{conjecture}{Conjecture}
\newtheorem{observation}{Observation}

\begin{document}

\title{Toward PK-Free Molecular Hamiltonians:\\
Two-Center Sturmian Integrals and Coupled Composition}
\thanks{Paper 19 in the Geometric Vacuum series}

\author{J.~Loutey}
\affiliation{Independent Researcher, Kent, Washington}

\date{April 3, 2026}

%% ========================================================================
%% ABSTRACT
%% ========================================================================

\begin{abstract}
The composed natural geometry framework (Paper~17) constructs
molecular Hamiltonians by partitioning electrons into single-center
blocks and coupling them via Phillips--Kleinman (PK) pseudopotentials.
This produces structurally sparse qubit Hamiltonians---$O(Q^{2.5})$
Pauli scaling, 51--1{,}712$\times$ fewer terms than Gaussian
baselines (Paper~14)---but PK limits accuracy to 5.3--26\%
equilibrium bond length error.  PK cannot be improved: six
modification attempts have failed (v2.0.6--23), and the overcounting
diverges at higher angular momentum.

We report a negative result (Track~CB, v2.0.37): replacing PK with
explicit cross-block electron-electron integrals \emph{without}
corresponding cross-center nuclear attraction integrals produces an
unbalanced Hamiltonian (29\% error, worse than PK's 15\%).  The root
cause is identified: two-center one-body integrals
$\langle \psi_{nlm}^A | V_B | \psi_{n'l'm'}^A \rangle$, where an
orbital on center~$A$ feels the nuclear potential of center~$B$, are
required for energy balance but absent from the single-center
framework.

We observe that these missing integrals are precisely the
Shibuya--Wulfman integrals of the Coulomb Sturmian
literature~\cite{shibuya1965,avery2004,avery2012}, which are computed
via Fock's momentum-space projection onto $S^3$---the same conformal
map that underlies the GeoVac framework (Paper~7).  The mathematical
foundation is therefore already in place.  We review the convergence
properties of the Shibuya--Wulfman expansion, identify the Coulomb
resolution technique~\cite{hoggan2011} as a potentially faster
alternative, and outline the ``coupled composition'' architecture
that would result from incorporating two-center integrals.  If the
angular expansion converges in few terms at molecular bond lengths,
the coupled Hamiltonian would be PK-free, structurally sparse, and
significantly more accurate than the current composed framework.
This paper documents the mathematical connections and research path;
implementation and benchmarks are deferred.
\end{abstract}

\maketitle

%% ========================================================================
\section{Introduction: The PK Problem}
\label{sec:intro}
%% ========================================================================

The composed natural geometry framework~\cite{loutey_paper17}
partitions electrons into groups---core, bond pairs, lone pairs---each
inhabiting its own natural coordinate system.  Inter-group coupling is
approximated by the Phillips--Kleinman (PK) pseudopotential, a
one-body operator that provides core-valence orthogonality and an
effective repulsive barrier.

PK is the framework's accuracy bottleneck.  At $l_{\max} = 2$ (the
optimal operating point), LiH achieves 5.3\% $R_{\rm eq}$ error;
BeH$_2$ gives 11.7\%; H$_2$O gives 26\%.  These errors are
structural, not numerical: the $l_{\max}$ divergence
(+0.15--0.22~bohr per additional angular channel) is confirmed to
originate in PK overcounting, not in the adiabatic approximation
(Track~BQ, v2.0.32).  Six attempts to modify PK have failed
(projector, spectral, self-consistent, $R$-dependent, $l$-dependent,
$l$-dependent on 2D solver), and the ab initio $R_{\rm ref}$
derivation failed because the relevant scale (${\sim}3$~bohr) is
molecular, not atomic (Track~AD, v2.0.18).

At the same time, PK is essential for the composed framework's
sparsity: it makes the electron repulsion integral (ERI) tensor
block-diagonal, yielding $O(Q^{2.5})$ Pauli scaling
(Paper~14)~\cite{loutey_paper14}.  Without PK, the block
decomposition has no coupling at all.

The question motivating this paper is: \emph{can PK be replaced by
explicit inter-block integrals that preserve structural sparsity while
eliminating the accuracy ceiling?}


%% ========================================================================
\section{The Negative Result: Coupled Composition Without $V_{ne}$}
\label{sec:negative}
%% ========================================================================

Track~CB (v2.0.37) tested the simplest version of this idea: remove
PK from the composed LiH Hamiltonian and add explicit cross-block
two-electron integrals $\langle pq | rs \rangle$ where $p, q$ belong
to one block and $r, s$ to another.

The cross-block ERIs were computed using existing infrastructure
(\texttt{cross\_block\_mp2.py}, originally built for Track~BX-3b).
Key census results for LiH at $n_{\max} = 2$, $Q = 30$:

\begin{table}[h]
\caption{Cross-block ERI census for LiH at $n_{\max} = 2$.}
\label{tab:census}
\begin{ruledtabular}
\begin{tabular}{lrr}
Metric & Within-block & Cross-block \\
\hline
Nonzero ERIs & 195 & 130 \\
Max $|\langle pq|rs\rangle|$ (Ha) & --- & 0.891 \\
Mean $|\langle pq|rs\rangle|$ (Ha) & --- & 0.142 \\
\end{tabular}
\end{ruledtabular}
\end{table}

The cross-block/within-block ERI ratio is 2/3, consistent with the
selection-rule counting bound of $\le 2\times$ from Paper~14,
Sec.~V.E.  Gaunt selection rules ($|\Delta l| \le 1$,
$|\Delta m| \le 1$) are preserved across block boundaries.

The resulting qubit Hamiltonian has 854 Pauli terms (2.56$\times$ the
composed value of 334) and 1-norm 85.7~Ha (2.30$\times$ the composed
value of 37.3~Ha).  However, the FCI energy tells the real story:

\begin{table}[h]
\caption{FCI accuracy for LiH at $n_{\max} = 2$ under different
coupling schemes.  Exact: $-8.071$~Ha.}
\label{tab:fci}
\begin{ruledtabular}
\begin{tabular}{lrr}
Configuration & $E_{\mathrm{FCI}}$ (Ha) & Error \\
\hline
No PK, no cross-block (decoupled) & $-7.195$ & 10.9\% \\
Composed (PK, no cross-block) & $-6.863$ & 15.0\% \\
Coupled (no PK, with cross-block $V_{ee}$) & $-5.734$ & 29.0\% \\
\end{tabular}
\end{ruledtabular}
\end{table}

The coupled Hamiltonian is \emph{worse} than both alternatives.  The
root cause is clear: cross-block ERIs add inter-group electron-electron
\emph{repulsion} without the balancing cross-center electron-nucleus
\emph{attraction}.  The Hamiltonian is energetically unbalanced.

\begin{observation}
\label{obs:decoupled}
The decoupled configuration (no PK, no cross-block interactions)
gives 10.9\% error---better than PK's 15.0\%.  This implies that
PK overcorrects at $n_{\max} = 2$: the blocks are more accurate
when left independent than when coupled via PK.
\end{observation}

This observation is consistent with the $l_{\max}$ divergence
diagnostic: PK's repulsive barrier systematically overcounts the
core-valence interaction energy, and the overcounting worsens with
basis size.


%% ========================================================================
\section{The Missing Piece: Two-Center Nuclear Attraction}
\label{sec:twocenter}
%% ========================================================================

The coupled Hamiltonian of Sec.~\ref{sec:negative} is incomplete.
A physically balanced molecular Hamiltonian requires, for each
electron:
\begin{enumerate}
\item Kinetic energy (within-block, already present).
\item Same-center nuclear attraction $V_A = -Z_A / r_A$ (within-block, already present).
\item Cross-center nuclear attraction $V_B = -Z_B / r_B$, where the
  electron on center~$A$ feels the potential of nucleus~$B$ (missing).
\item Electron-electron repulsion, both within-block (present) and
  cross-block (added in Track~CB).
\end{enumerate}

Item~3 requires two-center one-body integrals of the form
\begin{equation}
\label{eq:vne_cross}
I_{nlm, n'l'm'}^{AB} = \left\langle \psi_{nlm}^A \left|
  \frac{-Z_B}{|\mathbf{r} - \mathbf{R}_B|} \right|
  \psi_{n'l'm'}^A \right\rangle
\end{equation}
where $\psi_{nlm}^A$ is a hydrogenic orbital centered on nucleus~$A$,
and the potential is from nucleus~$B$ at position~$\mathbf{R}_B$.
These integrals are standard in the Slater-type orbital (STO)
literature but have never been computed in the GeoVac framework,
which has until now been strictly single-center at each level.


%% ========================================================================
\section{Connection to Coulomb Sturmian Theory}
\label{sec:sturmian}
%% ========================================================================

GeoVac's hydrogenic orbitals are Coulomb Sturmians with the
identification $k = Z/n$, where $k$ is the Sturmian exponent.
This is not merely a notational correspondence: the entire GeoVac
framework---the graph Laplacian on $S^3$ (Paper~7), the Fock
projection, the angular momentum selection rules---is built on the
same mathematical structure that the Avery group has developed for
molecular Coulomb Sturmian calculations since the
1990s~\cite{avery2000,avery2004,avery2012,avery2014}.

\subsection{Shibuya--Wulfman integrals}

When a Coulomb Sturmian centered on $A$ is expanded in the basis
centered on $B$, the expansion coefficients are the Shibuya--Wulfman
integrals~\cite{shibuya1965}:
\begin{equation}
\label{eq:sw}
S_{\tau', \tau} = \int d^3\mathbf{r}\; \chi_{\tau'}^*(\mathbf{r} -
\mathbf{R}_B)\, \frac{k}{r_A}\, \chi_\tau(\mathbf{r} - \mathbf{R}_A)
\end{equation}
where $\chi_\tau$ denotes a Coulomb Sturmian with quantum numbers
$\tau = (n, l, m)$, and the weighting $k/r_A$ arises from the
potential-weighted orthogonality of the Sturmian basis.

In Fock's momentum-space representation, Coulomb Sturmians map to
hyperspherical harmonics on $S^3$, and Shibuya--Wulfman integrals
become overlap integrals of translated hyperspherical harmonics.
The translation is parameterized by
\begin{equation}
S \equiv k \cdot R_{AB}
\end{equation}
where $R_{AB} = |\mathbf{R}_A - \mathbf{R}_B|$ is the internuclear
distance and $k$ is the common Sturmian exponent.

The nuclear attraction integral of Eq.~(\ref{eq:vne_cross}) is
closely related to the Shibuya--Wulfman integral: both involve a
Coulomb potential centered on a different nucleus acting on an orbital.
Koga and Matsuhashi~\cite{koga1987} derived sum rules for nuclear
attraction integrals over hydrogenic orbitals that express them as
finite sums over Clebsch--Gordan coefficients---the same angular
momentum algebra that generates Gaunt integrals in the GeoVac
framework.

\emph{Heterogeneous exponents.}  The standard Shibuya--Wulfman
formalism assumes a common Sturmian exponent~$k$ across all centers.
GeoVac's composed blocks use different effective charges
$Z_{\mathrm{eff}}$ per block, giving different exponents
$k_b = Z_{\mathrm{eff},b}/n$ for each block~$b$.  The generalized
overlap integrals between Sturmians of different exponents exist in
the literature---Koga and Matsuhashi~\cite{koga1987} treat this
case explicitly---but the convergence properties may differ from the
common-$k$ case.  Any implementation must use these generalized
integrals, and the convergence diagnostic (Sec.~\ref{sec:path})
must be performed with the actual heterogeneous exponents of the
LiH block decomposition, not with a uniform~$k$.

\subsection{The convergence challenge}

The Shibuya--Wulfman expansion has a well-documented convergence
limitation.  When a Sturmian on center~$A$ is expanded in the
angular basis of center~$B$, the expansion involves hyperspherical
harmonics of increasing angular momentum~$L$.  The rate of convergence
is controlled by the parameter $S = k \cdot R_{AB}$.

For $S \ll 1$ (nuclei close together relative to the orbital scale),
convergence is rapid.  For $S \gtrsim 3$ (typical molecular bond
lengths), convergence is slow, requiring many angular terms.  The
Avery group explicitly notes this: the expansion ``is very slowly
convergent if $S$ is large''~\cite{avery2012}.

For LiH at equilibrium ($R_{AB} \approx 3$~bohr), with $k \approx 1$
for the valence orbitals, $S \approx 3$.  This places the system
squarely in the slow-convergence regime.  Each additional angular
momentum channel in the expansion adds orbitals and therefore Pauli
terms to the qubit Hamiltonian.  If $L_{\max} = 5$--10 is required
for convergence, the Pauli count could approach the full Level~4N
encoding (${\sim}3{,}288$ terms), defeating the purpose of the block
decomposition.

\subsection{The Coulomb resolution alternative}

A potentially more efficient route was developed by Hoggan and
collaborators~\cite{hoggan2011}.  The ``Coulomb resolution'' technique
decomposes multi-center integrals over exponential-type orbitals into
sums of one-electron overlap-like products, each involving orbitals on
at most two centers.  Crucially, these two-center integrals are
separable in prolate spheroidal coordinates---the same coordinate
system used at Level~2 of the GeoVac hierarchy (Paper~11).

The Coulomb resolution sidesteps the slowly convergent angular
expansion of Shibuya--Wulfman by working directly in the two-center
coordinate system where the integrals are most naturally expressed.
GeoVac already possesses prolate spheroidal infrastructure from
Paper~11, including the Neumann expansion machinery for $V_{ee}$
(Paper~12).  The question is whether this infrastructure can be
extended to compute the one-body cross-center integrals of
Eq.~(\ref{eq:vne_cross}).


%% ========================================================================
\section{The Coupled Composition Architecture}
\label{sec:architecture}
%% ========================================================================

If two-center integrals can be computed efficiently, the resulting
``coupled composition'' Hamiltonian would have the following structure
for a molecule with $B$ blocks:

\begin{equation}
\label{eq:coupled}
H = \sum_{b=1}^{B} H_b^{(1)} + \sum_{b=1}^{B} V_{ee}^{(b)} +
\sum_{b < b'} V_{ee}^{(b,b')} + \sum_{b=1}^{B} \sum_{a \ne a_b}
V_{ne}^{(b,a)}
\end{equation}

where $H_b^{(1)}$ is the one-body Hamiltonian within block~$b$
(kinetic + same-center nuclear attraction), $V_{ee}^{(b)}$ is the
within-block ERI, $V_{ee}^{(b,b')}$ is the cross-block ERI, and
$V_{ne}^{(b,a)}$ is the cross-center nuclear attraction (electron in
block~$b$ feeling nucleus~$a$ from a different center).

Compared to the current composed Hamiltonian:
\begin{itemize}
\item PK is \emph{eliminated entirely}.  Core-valence orthogonality is
  enforced by the occupation number representation in second
  quantization, not by a pseudopotential.
\item Cross-block ERIs are included, subject to Gaunt selection rules.
  Track~CB established that these add ${\sim}2/3$ of the within-block
  ERI count (130 vs.~195 for LiH at $n_{\max} = 2$).
\item Cross-center nuclear attraction is included via two-center
  integrals.  The number of one-body terms is bounded by the number
  of orbital pairs times the number of off-center nuclei.
\end{itemize}

\begin{conjecture}
\label{conj:sparsity}
The Gaunt selection rules that enforce ERI sparsity within each block
also constrain the cross-block ERIs and the angular expansion of the
cross-center nuclear attraction.  If the two-center integral expansion
converges in $L_{\max} \le 3$ angular terms at molecular bond lengths,
the coupled Hamiltonian should have $\lesssim 1{,}500$ Pauli terms for
LiH at $n_{\max} = 2$---roughly $4.5\times$ the current composed
count of 334, but well below the full Level~4N count of 3{,}288 and
potentially below the Gaussian STO-3G count of 907 after the PK
one-body terms are removed.
\end{conjecture}

This conjecture is falsifiable by computing the two-center integrals
and performing the Jordan--Wigner transformation.

\begin{conjecture}
\label{conj:accuracy}
With a balanced Hamiltonian (all one-body and two-body terms present,
no PK), the FCI accuracy on the coupled Hamiltonian should be
significantly better than the composed framework's 5.3\% $R_{\rm eq}$
error, bounded below by the hydrogenic FCI accuracy floor
(${\sim}0.35\%$ for He at $n_{\max} = 5$, Paper~FCI-A).
\end{conjecture}


%% ========================================================================
\section{Relationship to Existing GeoVac Infrastructure}
\label{sec:infrastructure}
%% ========================================================================

The coupled composition program connects to multiple existing
components:

\emph{Paper~7 (Dimensionless Vacuum):} The $S^3$ conformal projection
that defines GeoVac's graph Laplacian is identical to Fock's
1935 momentum-space mapping~\cite{fock1935}, which is the foundation
of Shibuya--Wulfman theory.  The 18 symbolic proofs of Paper~7
validate the same mathematical structure that generates the
Shibuya--Wulfman integrals.

\emph{Paper~11 (Prolate Spheroidal):} The prolate spheroidal
coordinate system and Neumann expansion machinery are directly
applicable to the Coulomb resolution technique for two-center
integrals.

\emph{Paper~12 (Algebraic $V_{ee}$):} The Neumann expansion of
$1/r_{12}$ in prolate spheroidal coordinates demonstrates that
multi-center Coulomb integrals can be decomposed into algebraically
evaluable components via angular momentum recurrence.

\emph{Paper~14 (Qubit Encoding):} Section~V.E already bounds
cross-block ERI counts at $\le 2\times$ within-block; Track~CB's
empirical ratio of 2/3 is tighter.  The coupled Hamiltonian would
feed directly into the Jordan--Wigner encoding pipeline.

\emph{Paper~16 (Chemical Periodicity):} The $S_N$ representation
theory and atomic classifier determine the block decomposition.
Coupled composition uses the same block structure---it changes how
blocks are coupled, not how they are defined.

\emph{Paper~18 (Exchange Constants):} PK is classified as a
``calibration exchange constant''---a transcendental quantity that
enters when projecting between the graph and continuous coordinates.
If coupled composition eliminates PK, it removes one exchange constant
from the framework, consistent with the algebraicization program
(Sec.~XII of Paper~18).

\emph{Track~BX-3b (Cross-Block MP2):} The \texttt{cross\_block\_mp2.py}
module already computes cross-block ERIs for LiH.  This infrastructure
would be reused for the cross-block ERI component of the coupled
Hamiltonian.


%% ========================================================================
\section{Research Path}
\label{sec:path}
%% ========================================================================

The following steps are required to test the conjectures of
Sec.~\ref{sec:architecture}:

\begin{enumerate}
\item \textbf{Two-center integral implementation.}  Compute the
  nuclear attraction integrals of Eq.~(\ref{eq:vne_cross}) for
  hydrogenic orbitals.  Two approaches are available:
  (a)~direct Shibuya--Wulfman expansion via Clebsch--Gordan
  coefficients (well-documented but slowly convergent at $S \gtrsim 3$),
  or (b)~Coulomb resolution via prolate spheroidal decomposition
  (faster but requires new infrastructure building on Paper~11;
  note that Paper~12's Neumann expansion handles $V_{ee}$, not
  $V_{ne}$---the one-electron case is simpler but still requires
  implementation).
  \emph{Concrete diagnostic:} compute the Shibuya--Wulfman integrals
  for the LiH $1s$--$2s$ cross-center case at $R = 3.015$~bohr with
  the actual heterogeneous exponents ($k_{\mathrm{core}} = 3/1$,
  $k_{\mathrm{val}} = 1/2$), and measure how many $L$ channels are
  needed for 1\% relative convergence.
  \emph{Decision gate:} if $L_{\max} > 5$ is required for 1\%
  convergence on this simplest case, abort the direct SW route and
  proceed to Coulomb resolution (approach~b).

\item \textbf{Balanced coupled Hamiltonian.}  Construct the
  Hamiltonian of Eq.~(\ref{eq:coupled}) with all four interaction
  types present.  Perform Jordan--Wigner transformation and measure
  Pauli count, 1-norm, and QWC groups.

\item \textbf{FCI benchmark.}  Diagonalize the coupled Hamiltonian
  for LiH at $n_{\max} = 2$ ($Q = 30$, 4~electrons).  Compare energy
  and $R_{\rm eq}$ to exact, to composed (PK), and to Gaussian
  STO-3G.

\item \textbf{Scaling study.}  If Step~3 is positive, extend to
  $n_{\max} = 3$ and to BeH$_2$/H$_2$O to verify that the sparsity
  advantage persists and that accuracy improves for polyatomics.
\end{enumerate}

The critical risk is Step~1: if the angular expansion requires
$L_{\max} > 5$ for convergence, the Pauli count will approach the
full Level~4N encoding and the approach loses its structural advantage.
The Coulomb resolution route (approach~b) is the fallback if direct
expansion is too slow.


%% ========================================================================
\section{Discussion}
\label{sec:discussion}
%% ========================================================================

The composed framework's PK pseudopotential was a necessary compromise
for \emph{classical} solvers that could not handle the full
inter-group Coulomb interaction.  In second quantization on a quantum
computer, antisymmetry is enforced by the fermionic algebra, not by
coordinate-space projectors.  The occupation number representation
naturally prevents double-occupancy without PK.  This suggests that
PK's \emph{antisymmetry enforcement} role is a classical solver
artifact.  However, PK also approximates the physical core-valence
Coulomb interaction energy---this role must be \emph{replaced} by
explicit inter-block integrals, not simply removed.

Track~CB's negative result confirms this distinction.  The failure
was not because PK was needed for antisymmetry, but because the
replacement was incomplete: cross-block $V_{ee}$ was added without
the balancing cross-center $V_{ne}$.  The 10.9\% error of the fully
decoupled Hamiltonian (better than PK's 15.0\%) reinforces the
interpretation that PK overcorrects the Coulomb interaction, but the
remaining 10.9\% error shows that inter-block coupling \emph{is}
needed for quantitative accuracy.

A practical concern: Track~CB found that adding cross-block ERIs
increases the 1-norm from 37.3 to 85.7~Ha.  Even if PK is
eliminated, the net quantum resource cost could increase if the
balanced coupled Hamiltonian has higher 1-norm than the composed
Hamiltonian with PK classically partitioned (Track~BF: electronic-only
1-norm 33.3~Ha).  The sparsity advantage in Pauli count may be
partially offset by a 1-norm disadvantage.  This tradeoff must be
quantified for the balanced Hamiltonian.

The mathematical infrastructure for computing the missing integrals
exists in the Coulomb Sturmian literature and connects directly to
GeoVac's $S^3$ foundation.  The open question is convergence: does
the Shibuya--Wulfman expansion converge fast enough at molecular bond
lengths to preserve the structural sparsity that is GeoVac's primary
computational advantage?

If the answer is yes, the composed framework would be superseded by
coupled composition: same block structure, same Gaunt selection
rules, same $O(Q^{2.5})$ scaling class, but without PK and with
significantly better accuracy.  Paper~17's composed architecture
would then move to the archive, replaced by the coupled composition
as the production molecular Hamiltonian builder.

If the answer is no, the composed framework with PK remains the
production architecture, and the accuracy ceiling is accepted as the
structural cost of maintaining block-diagonal sparsity.  In either
case, this paper documents the mathematical connections and the
concrete data (Track~CB census, FCI benchmarks) that define the
research frontier.

\emph{Relationship to the LCAO negative result.}  The coupled
composition program reconstructs a multi-center molecular Hamiltonian
from atom-centered orbitals, which is superficially similar to the
LCAO approach abandoned in v0.9.x (Section~3 of the project
guidelines).  The key difference: LCAO used unstructured graph
concatenation with no natural geometry, producing $R$-independent
kinetic energies.  Coupled composition retains the natural geometry
block structure---each electron group inhabits its own optimal
coordinate system---and Gaunt selection rules enforce sparsity within
and across blocks.  The block decomposition provides physical
structure that LCAO lacked.


%% ========================================================================
%% BIBLIOGRAPHY
%% ========================================================================

\begin{thebibliography}{30}

\bibitem{loutey_paper7}
J.~Loutey, ``The Dimensionless Vacuum: Conformal Equivalence of
  Graph and Continuum Quantum Mechanics,'' GeoVac Paper~7 (2025).

\bibitem{loutey_paper11}
J.~Loutey, ``Prolate Spheroidal Lattice for Diatomic Molecules,''
  GeoVac Paper~11 (2026).

\bibitem{loutey_paper12}
J.~Loutey, ``Algebraic Two-Electron Integrals on the Prolate
  Spheroidal Lattice,'' GeoVac Paper~12 (2026).

\bibitem{loutey_paper14}
J.~Loutey, ``Structurally Sparse Qubit Hamiltonians from Spectral
  Graph Theory,'' GeoVac Paper~14 (2026).

\bibitem{loutey_paper17}
J.~Loutey, ``Composed Natural Geometries for Core-Valence Diatomics,''
  GeoVac Paper~17 (2026).

\bibitem{fock1935}
V.~Fock, ``Zur Theorie des Wasserstoffatoms,''
  Z.~Phys.\ \textbf{98}, 145 (1935).

\bibitem{shibuya1965}
T.~I.~Shibuya and C.~E.~Wulfman, ``Molecular orbitals in momentum
  space,'' Proc.\ Roy.\ Soc.\ A \textbf{286}, 376 (1965).

\bibitem{koga1987}
T.~Koga and T.~Matsuhashi, ``Sum rules for nuclear attraction
  integrals over hydrogenic orbitals,''
  J.~Chem.\ Phys.\ \textbf{87}, 4696 (1987).

\bibitem{avery2000}
J.~Avery, \emph{Hyperspherical Harmonics and Generalized Sturmians}
  (Kluwer Academic, Dordrecht, 2000).

\bibitem{avery2004}
J.~Avery, ``Many-center Coulomb Sturmians and Shibuya--Wulfman
  integrals,'' Int.\ J.\ Quantum Chem.\ \textbf{100}, 121 (2004).

\bibitem{hoggan2011}
P.~E.~Hoggan, ``How Exponential Type Orbitals Recently Became a Viable
  Basis Set Choice in Molecular Electronic Structure Work,'' in
  \emph{Solving the Schr\"odinger Equation: Has Everything Been
  Tried?}, edited by P.~Popelier (Imperial College Press, 2011),
  Chap.~7.

\bibitem{avery2012}
J.~S.~Avery and J.~E.~Avery, ``Coulomb Sturmians as a basis for
  molecular calculations,'' Mol.\ Phys.\ \textbf{110}, 1593 (2012).

\bibitem{avery2014}
J.~E.~Avery and J.~S.~Avery, ``Molecular integrals for Slater type
  orbitals using Coulomb Sturmians,''
  J.~Math.\ Chem.\ \textbf{52}, 301 (2014).

\bibitem{herbst2018}
M.~F.~Herbst, J.~E.~Avery, and A.~Dreuw, ``Quantum chemistry with
  Coulomb Sturmians: Construction and convergence of Coulomb Sturmian
  basis sets at Hartree--Fock level,'' arXiv:1811.05777 (2018).

\end{thebibliography}

\end{document}
